\setstretch{1.5}

Arabic optical character recognition (OCR) engines produce systematic errors
rooted in the visual properties of the Arabic script: dot-group character
confusions, morphologically valid but contextually wrong substitutions, and
segmentation failures that destroy word boundaries. Open-source engines such as
Qaari fall short of closed-source vision--language models, motivating the
question of whether post-OCR correction via a lightweight large language model
(LLM) can bridge this accuracy gap without modifying the recognition engine
and without task-specific fine-tuning.

This thesis presents a controlled, three-experiment multi-dataset evaluation of
knowledge-augmented prompt-engineering strategies for Arabic post-OCR correction
using Qwen3-4B-Instruct-2507 as the primary model. \textbf{Experiment~1}
evaluates eight prompt-engineering phases under an isolation principle across
nine datasets --- eight typewritten fonts from the PATS-A01 corpus (Qaari
baseline CER 2.12\%--12.12\%) and one handwritten corpus (KHATT, 34.24\%
baseline CER) --- each changing exactly one variable relative to a zero-shot
baseline. The five augmentation conditions are: (1)~failure-filtered character
confusion matrix injection, cross-referenced with the model's zero-shot
failures (Phase~3); (2)~self-reflective prompting derived from training-split
error statistics including quantified fix and introduction rates, concrete
word-level error pairs, and an explicit overcorrection warning (Phase~4);
(3)~CAMeL Tools morphological post-processing with position-indexed
known-overcorrection revert (Phase~5); (4)~BM25 character $n$-gram
retrieval-augmented generation using same-domain training corrections (Phase~8);
and (5)~error-signature RAG combining CAMeL profile and confusion-matrix
character similarity (Phase~9). A combination and ablation study (Phase~6)
tests four multi-component configurations. \textbf{Experiment~2} extends
evaluation to thirteen datasets, adding four full-page domain-diverse corpora:
KHATT-Paragraph (449 handwritten paragraph images), Yarmouk (1{,}663 contemporary
printed pages), Muharaf (116 handwritten manuscript images), and Historical
(40 archival manuscript pages from eight Arabic books). Eight prompt
configurations are tested in a factorial design. \textbf{Experiment~3}
benchmarks four LLMs (Qwen3-4B, Qwen3-14B, Gemma-3-4B, Gemma-3-12B) across
two OCR sources (Qaari, Gemma-3 VLM) and two held-out benchmarks
(RDI-Test-Lines, Kitab) to assess scalability and domain generalisation.

All eight PATS-A01 fonts share 2,766 text lines drawn from the same underlying
content. An aligned 85/15 train/validation split (seed 42) ensures that every
font receives the same 415 validation samples, so that font is the only
variable when comparing OCR quality and correction effectiveness across
typefaces. KHATT provides a complementary handwritten evaluation with
substantially higher baseline error rates.

All metrics use the primary evaluation protocol: normal samples only
(OCR/GT length ratio $\leq 5.0$, excluding Qaari repetition-runaway samples),
with diacritics stripped before CER/WER computation.

Phase~1 establishes the OCR error landscape. Baseline CER on PATS-A01 fonts
ranges from 2.12\% (Akhbar) to 12.12\% (Andalus), with a macro-average of
5.57\% (WER 13.78\%). KHATT handwritten text reaches 34.24\% CER (WER 75.60\%).
The dominant error type across PATS-A01 is dot confusion (42.1\% of
substitutions), followed by similar-shape confusions (18.3\%) and hamza
errors (11.6\%). Morphological analysis reveals that 71.3\% of OCR errors
produce valid Arabic words in the wrong context (\textit{valid-but-wrong}),
with only 28.7\% producing morphologically invalid tokens.

Phase~2 (zero-shot correction) \textit{increases} the PATS-A01 macro-average
CER from 5.57\% to 5.84\% ($+$4.8\% relative) --- a net-harmful outcome on
average. The effect is strongly font-dependent: on Traditional font, CER drops
from 3.63\% to 2.96\% ($-$18.5\% relative), while on Akhbar (2.12\% baseline)
zero-shot correction increases CER to 2.71\%. On KHATT, correction is
beneficial: 34.24\%~$\to$~33.25\% ($-$0.99~pp).

Phase~3 (confusion-matrix injection) achieves PATS-A01 macro-average 5.74\%
($-$0.10~pp vs.\ Phase~2) and 33.26\% on KHATT. Phase~4 (self-reflective
prompting) achieves the strongest isolated Experiment~1 result: 5.59\% on
PATS-A01 ($-$0.25~pp vs.\ Phase~2, near-parity with the 5.57\% OCR baseline)
and 33.24\% on KHATT. Phase~5 (CAMeL morphological post-processing) results
were invalidated by implementation bugs and are excluded. Phase~6 combinations
show that Phase~3 and Phase~4 address orthogonal failure modes and combine
near-additively. The \texttt{conf\_self} combination achieves 5.76\% PATS and
33.27\% KHATT. Phase~8 (BM25 RAG) achieves the best Experiment~1 PATS result:
5.45\% ($-$0.12~pp below the OCR baseline), the \emph{only} Experiment~1
strategy to beat the OCR baseline. Phase~9 (Error-Signature RAG) achieves the
best KHATT result: 33.13\% ($-$1.11~pp vs.\ OCR baseline).

The central finding is the identification and formal characterisation of an
\emph{over-correction threshold}: a baseline CER level ($e^*$) below which the
LLM introduces more errors than it corrects. The threshold arises because the
introduction rate (fraction of correct tokens incorrectly changed) is
approximately constant across fonts, while the fix rate scales with error
density. Empirically, $e^*$ lies at approximately 6--7\% CER for this
model--engine pair: Akhbar (4.3\% baseline) is harmed by every strategy
tested, while Simplified (7.9\%) benefits. Phase~4 is the only strategy that
partially mitigates over-correction on Akhbar, because its diagnostic text
explicitly targets the introduction rate rather than only providing correction
guidance.

The ablation study reveals that Phases~3 and~4 address orthogonal failure
modes---Phase~3 targets the OCR engine's known confusion patterns (raising the
fix rate), while Phase~4 targets the correction model's systematic biases
(lowering the introduction rate)---and combine near-additively. On KHATT
handwritten text, where 37\% of errors are segmentation failures (word merges
and splits that destroy word boundaries), text-only correction is structurally
limited regardless of prompt strategy, achieving only 3.6--5.0\% relative
improvement despite a 44.6\% baseline CER.

Experiment~2 (eight-trial prompt design study across all thirteen datasets)
finds that conservative zero-shot prompting (T3) achieves the best PATS-A01
result: 5.27\% CER ($-$5.4\% relative vs.\ OCR baseline), outperforming all
Experiment~1 phases without retrieval overhead. The base+RAG configuration
(T2) is the best overall (5.44\% PATS, 33.82\% KHATT) and is selected for
Experiment~3. On the four full-page datasets, T2 reduces the macro-average
from 86.50\% to 68.18\% ($-$21.2\% absolute) after runaway correction.
Error-pattern trials (T4--T8) are brittle: T7 fails catastrophically
(190.60\% PATS CER, prompt-template injection).

Experiment~3 (four LLMs, two OCR sources, two held-out benchmarks) finds
that Gemma-3-12B achieves the best PATS CER$\dagger$ (5.33\%); Qwen3-4B
is best on KHATT (33.82\%) and full-page datasets (68.18\% vs.\ OCR
baseline 86.50\%). Qwen3-14B suffers a runaway epidemic on KHATT (175.69\%
raw; 35.55\% after correction). LLM correction generalises robustly to
unseen benchmarks: RDI-Test-Lines $-$10.7\% absolute, Kitab Benchmark
$-$17\% to $-$24\% relative. Gemma VLM OCR outperforms Qaari on
RDI-Test-Lines (68.49\% vs.\ 95.31\%), demonstrating stronger generalisation
on diverse document styles.

This thesis contributes: (1)~a formal over-correction threshold model with
empirical estimation ($e^*\approx6$--$7\%$ CER for Qwen3-4B on Qaari);
(2)~a novel self-reflective prompting strategy achieving the strongest
Experiment~1 isolated result ($-$0.25~pp PATS, only strategy to improve
Akhbar); (3)~a proof of complementarity between confusion-matrix and
self-reflective augmentation (orthogonal failure modes, near-additive
contributions); (4)~a comprehensive thirteen-dataset evaluation spanning
typewritten, handwritten, historical manuscript, encyclopedic printed, and
paragraph-level domains across three experiments; (5)~a three-stage
export--infer--analyse pipeline enabling GPU-free local reproducibility;
(6)~the finding that conservative prompting outperforms knowledge augmentation
for clean typewritten text (T3: 5.27\% PATS, $-$5.4\% relative, no retrieval
overhead); and (7)~an empirical characterisation of LLM correction
generalisation to unseen benchmarks and multiple model scales.
